\section{Positividades, espacios físicos y transporte del espectro}
La geometría positiva y la teoría gauge reciben estructuras distintas de
un mismo proceso genealógico. La primera realiza los pesos y sus
incidencias mediante menores ordenados, celdas y formas canónicas.
La segunda conserva la incidencia gauge, la holonomía, las
representaciones, las rutas y la memoria dentro de una medida compatible.
Su forma de positividad relevante para la reconstrucción física es la
positividad por reflexión. El nexo conceptual adquiere contenido
matemático cuando se conservan los mapas que producen cada forma.

En la parte geométrica se ha demostrado que la subdivisión de una celda
conserva su forma canónica mediante cancelación de polos interiores.
Para las energías nodales, la reducción de Schur conserva una forma
efectiva y registra el detalle. En la parte gauge, la compatibilidad del
refinamiento debe conservar además el vacío, el espacio de invariantes,
el generador y el testigo espectral. Ésta es la cadena desarrollada a
continuación: escala aritmético--geométrica, medida común, semigrupo,
proyección física, refinamiento y reconstrucción continua.

La norma de un vector, la positividad de una matriz de momentos y la
existencia de una brecha de un Hamiltoniano poseen dominios diferentes.
Una congruencia con factor invertible conserva la positividad y la
inercia de una forma; un entrelazamiento isométrico del generador, con su
dominio, transporta su realización sobre la imagen reducida de la
inclusión. La cota uniforme y el testigo persistente fijan después la
brecha límite. Por ello el desarrollo mantiene visibles las
inclusiones y las identidades de semigrupo que justifican el paso al
límite. La escala de la brecha procede de la incidencia areal--volumétrica
y de su unidad dimensional declarada.

Los cuerpos siguientes reúnen la construcción proyectiva y su
desarrollo posterior dependiente del acoplamiento \(g\). Las primeras
comparaciones con una parametrización Gibbs se completan después mediante
la acción, la medida y el transporte explícitos. El orden expositivo
conserva las demostraciones del proceso y sitúa la formulación completa
tras sus dependencias. Las cuatro reflexiones corresponden a las cuatro
direcciones euclídeas de la reconstrucción. El desarrollo aquí reunido
tiene por objeto esa teoría gauge y su brecha, con el grupo compacto,
conexo y simple declarado en los teoremas.

El sector supersimétrico planar en el que se estudian amplitudes mediante
amplituedros conserva una pregunta física específica. Para enlazarlo con
la presente realización se requieren los operadores de supercarga, sus
dominios y la representación correspondiente. La dimensión cuatro de la
coordenatización euclídea y la diferencia cuatro del bloque modular permanecen
resultados ya definidos de sus respectivos operadores.
